\begin{helper}[Original deleted-block avoidance]
Let $v\in Z$, let $\beta$ be missing at $v$, free for
$(T,e_0,\varphi)$, and distinct from $\gamma(\delta)$, and let $P$ be a
full simple $(\delta,\beta)$ path from $v$ to an outside endpoint.
Neither vertex of the deleted block $U_{\gamma(\delta)}$ occurs in the
support of $P$.
Consequently no $\delta$-colored boundary edge of $V(T)$ is incident
with either vertex of that block.  Both conclusions retain exact vertex
and parallel-edge identities.
\end{helper}
\begin{proof}
Put
\[
 X=\noderef{TreeSequenceVertexSet}(T),\qquad
 U=\{u^1_{\gamma(\delta)},u^2_{\gamma(\delta)}\},\qquad
 Z=X\setminus U.
\]
The hypothesis says that $v\in Z$, that $\beta$ is missing at $v$ and
free for $T$, and that $\beta\ne\gamma(\delta)$.

We first prove carefully that no $\beta$-edge joins $U$ to $Z$.  By
\noderef{ScheideStructuralInternalPair}, there is an individual edge
identity $b$ such that
\[
 \varphi(b)=\beta,
 \qquad
 G.\operatorname{ends}(b)
   =\{u^1_{\gamma(\delta)},u^2_{\gamma(\delta)}\}.       \tag{1}
\]
Unpack the supplied \noderef{ScheideReducedConfiguration}.  As a reduced
structural configuration it contains the underlying
\noderef{ScheideStructuralConfiguration}, and in particular its explicitly
named original globally maximum triple $(T,e_0,\varphi)$.  Unpacking
\noderef{GlobalMaximumTriple} gives the original $k$-triple
$(G,e_0,\varphi)$.  Unpacking \noderef{KTriple} in turn gives that
$\varphi$ is proper in the sense of
\noderef{IsProperEdgeColoring}.

Fix $u\in U$.  Equation (1) makes $b$ incident with $u$ in the sense of
\noderef{Incident}.  If a distinct edge identity $a$ of color $\beta$
were also incident with $u$, then $a$ and $b$ would be adjacent by
\noderef{AdjacentEdges}.  This contradicts the just-extracted properness.
Thus $b$ is the unique $\beta$-edge incident with either vertex of $U$.
Since the two exact ends of $b$ are the two vertices of $U$, no
$\beta$-edge can join $U$ to $Z$.  Independently,
\noderef{ScheideDefectiveNoCross}, instantiated with the present defective
color $\delta$, says that no $\delta$-edge has one incident vertex in
$U$ and another in $Z$.  Hence neither of the two path colors crosses
directly between $Z$ and $U$.                                      \tag{2}

Apply \noderef{ScheideSymmetricOriginalPathOrder} to the present
$\beta,v,z,P$.  It supplies exact path indices for the consecutive
five-node excursion
\[
 f^1_\delta,u^1_\delta,h,u^2_\delta,f^2_\delta
\]
or its reverse, where $h$ is the internal $\beta$-edge of $U_\delta$.
It also supplies exact indices $q_v,q_f$ such that $q_f=q_v+1$, the
node at $q_v$ is $v^0_\delta$, the node at $q_f$ is $f^0_\delta$, the
excursion ends before $q_v$, and every vertex-node strictly after $q_v$
lies outside $X$.                                                   \tag{3}

Suppose for a contradiction that a vertex of $U$ occurs on $P$.  Among
its finitely many path indices choose the first such index $t$.  Since
the initial vertex $v$ lies in $Z$, there is a path index before $t$
whose node is a vertex of $Z$; choose the last such index $q$.  By the
bipartite adjacency definition in
\noderef{TwoColorIncidenceGraph}, the node at $q+1$ is an edge-node,
say the identity $f$, and the node at $q+2$ is a vertex-node, say $x$.
The two walk adjacencies say that $f$ is incident with both the $Z$-vertex
at $q$ and $x$, and that its color is $\delta$ or $\beta$.  Simplicity
of the supplied full path and the endpoint-pair convention in
\noderef{LooplessMultigraph} identify these as the two exact ends of
$f$.

The vertex $x$ lies outside $X$.  Indeed, if $x\in Z$, then its index
$q+2\le t$ is in fact strictly below $t$ and contradicts the maximal
choice of $q$.  If instead $x\in U$, then $f$ crosses directly from $Z$
to $U$, contradicting (2), whichever of its two possible colors it has.
Since $X=Z\mathbin{\dot\cup}U$, these are all possibilities inside $X$.
Thus the definition \noderef{BoundaryEdges} makes $f$ a boundary edge of
both $X$ and $Z$.                                                   \tag{4}

Instantiate \noderef{ScheideGlobalMaximumOrder} with the original
globally maximum triple extracted above.  Its last clause says that $X$
is closed.  By \noderef{ColoredVertexSetClosed}, because $\beta$ is
missing at $v\in X$, this closedness forbids a $\beta$-colored boundary
edge of $X$.  Consequently (4) and the two-color incidence definition
force $\varphi(f)=\delta$.  The
exact-three clause of \noderef{ScheideReplacementData}, applied to the
$Z$-end and the other end outside $X$, now gives
\[
                 f\in\{f^0_\delta,f^1_\delta,f^2_\delta\}.          \tag{5}
\]

Suppose first that $f$ is $f^1_\delta$ or $f^2_\delta$.  Its occurrence
at $q+1$ cannot be the second edge of the excursion in (3): the node
immediately before that second edge is one of the outside vertices
$u^1_\delta,u^2_\delta$, whereas the node at $q$ lies in $Z\subseteq X$.
Thus it is the first excursion edge.  The exact consecutive indices in
(3) carry the path through the two outside vertices and the internal
$\beta$-edge to the other named $\delta$-edge.  That other edge has exact
ends $\{v^i_\delta,u^i_\delta\}$ by
\noderef{ScheideReplacementData}; simplicity prevents the next node from
returning to the immediately preceding outside vertex, so the next node
is $v^i_\delta\in Z$.  It occurs after $q$ and still before $t$, contrary
to the maximal choice of $q$.

The remaining alternative in (5) is $f=f^0_\delta$.  Simplicity gives
this edge-node a unique path index, so (3) identifies $q+1$ with $q_f$.
Therefore $t>q_f>q_v$.  But the node at $t$ is a vertex of $U\subseteq X$,
whereas the last clause of (3) says that every vertex-node after $q_v$
is outside $X$.  This is again a contradiction.  All alternatives in
(5) are impossible, proving
\[
       u^1_{\gamma(\delta)},u^2_{\gamma(\delta)}\notin V(P).        \tag{6}
\]

Finally instantiate \noderef{ScheideDefectiveChainBoundaryEdges} with
the same explicitly identified original globally maximum triple, the
defective color $\delta$, the free color $\beta$, its missing witness
$v$, and the supplied full path.  It says that every individual
$\delta$-colored boundary identity of $X$ occurs as an edge-node in the
support of $P$.  If such an identity $a$ were incident with a vertex
$u\in U$, then \noderef{TwoColorIncidenceGraph} would make the
vertex-node $u$ adjacent to the edge-node $a$.  Hence they lie in the
same two-color component.  The support-equality clause in
\noderef{ScheideDefectiveChainPath} identifies that entire component
with the support of $P$, so the vertex-node $u$ would lie on $P$.
This contradicts (6) and proves the required boundary avoidance, with
the individual identity $a$ retained throughout.
\end{proof}
